% ----------------- Paquetes -----------------
\usepackage[utf8]{inputenc}
\usepackage[spanish]{babel}
\usepackage[a4paper,top=2.5cm,bottom=2.5cm,left=3cm,right=3cm]{geometry}
\usepackage{setspace}
\usepackage{graphicx}
\usepackage{amsmath}
\usepackage{titlesec}
\usepackage{url}
\usepackage{enumitem}
\usepackage{multicol}
\usepackage{helvet} % Arial-like font
\usepackage{comment}
\usepackage{caption}
\usepackage{indentfirst}
% ----------------- Fuente Arial -----------------
\usepackage{fancyhdr} % Encabezado y pie de página
\renewcommand{\familydefault}{\sfdefault}
% ----------------- Fuente Arial -----------------

% ----------------- Interlineado 1.5 -----------------
\onehalfspacing

% ----------------- Espaciado de párrafo (0 pt antes/después) -----------------
\setlength{\parskip}{0pt}

% ----------------- Sangría de 0.5 cm al inicio de sección/subsección -----------------
\setlength{\parindent}{0.5cm}

% ----------------- Estilos de sección -----------------
\titleformat{\section}{\bfseries\fontsize{14}{16}\selectfont}{\thesection}{1em}{}
\titleformat{\subsection}{\bfseries\fontsize{12}{14}\selectfont}{\thesubsection}{1em}{}

% ----------------- Nombre de Tabla en español -----------------
\addto\captionsspanish{\renewcommand{\tablename}{Tabla}}

% ----------------- Formato de títulos de figuras y tablas -----------------
% Arial 11, sin dos puntos, centrado. El punto final se escribe a mano
% en cada \caption{...}: con format=puntofinal el punto se colocaba fuera
% del bloque centrado del caption y aparecía huérfano en su propia línea.
\DeclareCaptionFont{arial11}{\fontsize{11}{13}\selectfont}
\captionsetup{
  font=arial11,
  labelfont=arial11,
  labelsep=space,
  justification=centering,
  singlelinecheck=false
}

% ----------------- Nota de fuente (Arial 9, cursiva, alineada a la derecha) -----------------
\newcommand{\fuente}[1]{\par\vspace{2pt}{\setstretch{1.15}\raggedleft\fontsize{9}{11}\selectfont\itshape Fuente: #1\par}\vspace{4pt}}

% ----------------- Pie de página -----------------
\pagestyle{fancy}
\fancyhf{} % Limpia encabezados y pies
\fancyfoot[L]{\fontsize{10}{12}\selectfont Circuitos Eléctricos de CA, IE-AD25}
\fancyfoot[R]{\fontsize{10}{12}\selectfont Dr. Luis Alejandro Flores Oropeza}
\renewcommand{\headrulewidth}{0pt} % sin línea en encabezado
\renewcommand{\footrulewidth}{0pt} % sin línea en pie de página

%-- Das comm --
\newcommand{\titulos}[3]{
  {\fontsize{12}{14}\bfseries \textit{#1}}\\
  {\fontsize{10}{12}\normalfont #2}\\
  {\fontsize{9}{11}\itshape #3}\\[0.5cm]

}